\chapter{Fourier Series}\label{ch:fourier}

\begin{center}\itshape
After $1000$ terms, $\sin 1+\frac{\sin 2}2+\frac{\sin 3}3+\cdots$ equals $1.0707$.\\
The sums are circling $(\pi-1)/2$, which equals $1.0708$.
\end{center}

Add the terms yourself. Angles are in radians. Ten terms give $1.1245$. A hundred terms give $1.0604$. A thousand terms give $1.0707$. The partial sums wander, slowly, around $1.0708$. Start again with the even angles, $\sin 2+\frac{\sin 4}2+\frac{\sin 6}3+\cdots$. Those sums settle near $0.5708$, which is $\frac{\pi-2}2$. A sum of smooth waves is drawing the straight line $\frac{\pi-x}2$. Example~\ref{ex:sawtooth} proves the identity, and it shows where the line has to break.

In 1807 Joseph Fourier told the Paris Academy that a function on an interval is a sum of sines and cosines, even when different formulas describe different pieces. Lagrange objected. Fourier could not yet say in which sense the sum converges. The question occupied analysts for a century. It forced the modern definitions of function, of integral, and of uniform convergence, and it led Cantor to set theory. In the preliminary discourse of \emph{Th\'eorie analytique de la chaleur} (1822), Fourier wrote that the profound study of nature is the most fertile source of mathematical discoveries.

\portrait{fourier}{Joseph Fourier (1768--1830). Engraving by Am\'ed\'ee Geille, after Julien-L\'eopold Boilly. Public domain.}

Three theorems answer Fourier's question, in three senses that a computation can use. At a point, the series of a piecewise $C^1$ function converges to the function, and to the average of the two limits at a jump. That is Dirichlet's theorem. Uniformly, the averages of the partial sums of a continuous function converge to the function. That is Fej\'er's theorem. The partial sums themselves may still diverge. In mean square, the series of a piecewise continuous function converges to it, and the energy of the function is the sum of the energies of its harmonics. That is Parseval's identity. The same orthogonality gives the Gibbs overshoot, Weyl's test for equidistribution, the return probability of a walk on $\Z^d$ promised in Volume~4, and Hurwitz's proof that the circle encloses the largest area for a given length.

\begin{think}
The functions $\sin x,\sin 2x,\sin 3x,\dots$ are all smooth. Can a sum of them have a jump? If it has one, what value do you expect the series to take at the jump?
\end{think}

Yes, and the value is the average of the two limits. The opening sum is $0$ at multiples of $2\pi$, where the sawtooth jumps, and it equals $\frac{\pi-x}2$ between the jumps. The bump beside the jump does not shrink. Its height is about $1.179$ on the square wave, for every $N$.

\section{Fourier coefficients and best approximation}

Throughout, functions are $2\pi$-periodic, from $\R$ to $\C$. Such a function is \emph{piecewise continuous} when $[-\pi,\pi]$ splits into finitely many intervals on whose interiors $f$ is continuous, with finite limits at both ends. Its integral is the sum of the integrals over the pieces. It is \emph{piecewise $C^1$} when, on each piece, $f$ coincides with a $C^1$ function on the closed interval. Since $f$ is periodic, its integral over any interval of length $2\pi$ is the same. Cut the interval at a multiple of $\pi$ and translate one part by $2\pi$.

For such functions let
\[
\begin{gathered}\langle f,g\rangle=\frac1{2\pi}\int_{-\pi}^\pi f(t)\overline{g(t)}\dd t,\qq\|f\|_2=\sqrt{\langle f,f\rangle},\end{gathered}
\]
and let $e_n(t)=e^{int}$ for $n\in\Z$. The integral of $e^{ikt}$ over $[-\pi,\pi]$ is $0$ for every integer $k\ne0$, so the $e_n$ are \emph{orthonormal}: $\langle e_m,e_n\rangle$ is $1$ if $m=n$ and $0$ otherwise. The form $\langle\cdot,\cdot\rangle$ is linear in its first argument and Hermitian, and $\langle f,f\rangle\ge0$. Cauchy--Schwarz holds, with the proof from Volume~1, Chapter~9: $|\langle f,g\rangle|\le\|f\|_2\|g\|_2$. A piecewise continuous $f$ can have $\|f\|_2=0$ without vanishing, after a change at finitely many points. That does no harm below.

\begin{definition}
The \emph{Fourier coefficients} of $f$ are
\[
c_n(f)=\langle f,e_n\rangle=\frac1{2\pi}\int_{-\pi}^\pi f(t)e^{-int}\dd t,
\]
and its \emph{partial Fourier sums}
are $S_Nf=\sum_{|n|\le N}c_n(f)e_n$. A \emph{trigonometric polynomial of degree at most $N$} is a combination of $e_{-N},\dots,e_N$.
\end{definition}

Try $f(x)=\cos x$. Then $\cos x=\frac12 e^{ix}+\frac12 e^{-ix}$, so $c_1(f)=c_{-1}(f)=\frac12$ and every other coefficient is $0$. Hence $S_Nf=\cos x$ for every $N\ge1$. The function is already its own best approximation of degree $1$.

For real $f$, grouping $n$ with $-n$ gives the real form $S_Nf(x)=\frac{a_0}2+\sum_{n=1}^N(a_n\cos nx+b_n\sin nx)$, with
$a_n=\frac1\pi\int f(t)\cos nt\dd t$ and $b_n=\frac1\pi\int f(t)\sin nt\dd t$.

\begin{proposition}[best approximation]\label{prop:bestapprox}
Let $f$ be piecewise continuous and $P$ a trigonometric polynomial of degree at most $N$. Then
\[
\|f-P\|_2^2=\|f-S_Nf\|_2^2+\|S_Nf-P\|_2^2.
\]
In particular $S_Nf$ is the trigonometric polynomial of degree at most $N$ closest to $f$, and
$\|f-S_Nf\|_2^2=\|f\|_2^2-\sum_{|n|\le N}|c_n(f)|^2$.
\end{proposition}
\begin{proof}
For $|n|\le N$, orthonormality gives $\langle f-S_Nf,e_n\rangle=c_n(f)-c_n(f)=0$. So $u=f-S_Nf$ is orthogonal to $v=S_Nf-P$, and $v$ is a combination of these $e_n$. Then
\[
\|u+v\|_2^2=\|u\|_2^2+\|v\|_2^2+2\operatorname{Re}\langle u,v\rangle=\|u\|_2^2+\|v\|_2^2.
\]
With $P=0$ this reads $\|f\|_2^2=\|f-S_Nf\|_2^2+\|S_Nf\|_2^2$, and orthonormality gives $\|S_Nf\|_2^2=\sum_{|n|\le N}|c_n(f)|^2$.
\end{proof}

The identity is Pythagoras. The error $f-S_Nf$ is perpendicular to every trigonometric polynomial of degree at most $N$, so moving from $S_Nf$ to any other such polynomial only adds length.

\begin{corollary}[Bessel's inequality and Riemann--Lebesgue]\label{cor:bessel}
For $f$ piecewise continuous, $\sum_{n\in\Z}|c_n(f)|^2\le\|f\|_2^2$. Hence $c_n(f)\to0$ as $|n|\to\infty$, and also
\[
\int_{-\pi}^\pi f(t)\sin\Bigl(N+\frac12\Bigr)t\dd t\xrightarrow[N\to\infty]{}0.
\]
\end{corollary}
\begin{proof}
The partial sums of $\sum|c_n|^2$ are at most $\|f\|_2^2$, by Proposition~\ref{prop:bestapprox}. The terms of a convergent series tend to $0$. For the last claim, expand
\[
\sin\Bigl(N+\frac12\Bigr)t=\sin Nt\cos\frac t2+\cos Nt\sin\frac t2.
\]
The functions $\sin Nt$ and $\cos Nt$ are combinations of $e_{\pm N}$, so the integral is a combination of $c_{\pm N}\bigl(f\cos\frac t2\bigr)$ and $c_{\pm N}\bigl(f\sin\frac t2\bigr)$, with bounded coefficients. Those coefficients tend to $0$.
\end{proof}

This completes the Riemann--Lebesgue lemma, which was an exercise of Volume~1, Chapter~8, for $f$ of class $C^1$. Here $f$ is only piecewise continuous. The proof is Pythagoras in a space of functions. Volume~1 reached the same limit by an integration by parts, for a smoother $f$.

\section{Dirichlet's theorem}

The partial sum $S_Nf(x)$ is an average of the values of $f$ against a kernel that depends only on the difference of the two variables. The kernel has a closed form, and that form is what makes the jump computable.

\begin{lemma}[Dirichlet kernel]\label{lem:dirichletkernel}
For every $N\ge0$ and every $x$, $S_Nf(x)=\frac1{2\pi}\int_{-\pi}^\pi f(x+u)D_N(u)\dd u$, where
\[
D_N(u)=\sum_{|n|\le N}e^{inu}=\frac{\sin(N+\frac12)u}{\sin\frac u2}
\]
for $u\notin2\pi\Z$, and $D_N(0)=2N+1$. Moreover $D_N$ is even and $\frac1{2\pi}\int_0^\pi D_N(u)\dd u=\frac12$.
\end{lemma}
\begin{proof}
By definition,
\[
\begin{aligned}S_Nf(x)&=\frac1{2\pi}\int_{-\pi}^\pi f(t)\sum_{|n|\le N}e^{in(x-t)}\dd t\eqbr=\frac1{2\pi}\int_{-\pi}^\pi f(t)D_N(t-x)\dd t,\end{aligned}
\]
and $D_N$ is even. Substitute $t=x+u$ and use periodicity. For the closed form, multiply the sum by $e^{iu/2}-e^{-iu/2}$. The product telescopes:
\[
D_N(u)\bigl(e^{iu/2}-e^{-iu/2}\bigr)=e^{i(N+\frac12)u}-e^{-i(N+\frac12)u}.
\]
The left side is $2i\sin\frac u2\,D_N(u)$ and the right side is $2i\sin\bigl(N+\frac12\bigr)u$, so $D_N(u)=\frac{\sin(N+\frac12)u}{\sin\frac u2}$. Finally $\frac1{2\pi}\int_{-\pi}^\pi D_N=1$, since only the term $n=0$ contributes, and the two halves are equal because $D_N$ is even.
\end{proof}

\begin{theorem}[Dirichlet, 1829]\label{thm:dirichletfourier}
Let $f$ be piecewise $C^1$. For every $x$,
\[
S_Nf(x)\xrightarrow[N\to\infty]{}\frac{f(x^+)+f(x^-)}2.
\]
In particular $S_Nf(x)\to f(x)$ wherever $f$ is continuous.
\end{theorem}
\begin{proof}
\emph{Setup.} By Lemma~\ref{lem:dirichletkernel}, and writing $\frac12f(x^\pm)=\frac1{2\pi}\int f(x^\pm)D_N$ over $[0,\pi]$ and $[-\pi,0]$ respectively,
\[
\begin{aligned}S_Nf(x)&-\frac{f(x^+)+f(x^-)}2\eqbr=\frac1{2\pi}\int_{-\pi}^\pi G(u)\sin\Bigl(N+\frac12\Bigr)u\dd u,\end{aligned}
\]
where $G(u)=\frac{f(x+u)-f(x^+)}{\sin(u/2)}$ for $0<u\le\pi$ and $G(u)=\frac{f(x+u)-f(x^-)}{\sin(u/2)}$ for $-\pi\le u<0$.

\emph{The quotient.} Away from $0$, the function $G$ is piecewise continuous. Near $0^+$, the function $f$ coincides on some $\left]x,x+\eta\right]$ with a $C^1$ function $h$ on $[x,x+\eta]$ with $h(x)=f(x^+)$, so
\[
G(u)=\frac{h(x+u)-h(x)}u\cdot\frac u{\sin(u/2)},
\]
which tends to $2h'(x)$ as $u\to0^+$. The same limit holds at $0^-$. So $G$ extends to a piecewise continuous function on $[-\pi,\pi]$.

\emph{Finish.} The integral tends to $0$ by Corollary~\ref{cor:bessel}.
\end{proof}

Convergence at $x$ depends only on the behaviour of $f$ near $x$. This is the \emph{localisation principle}. Far from $u=0$, the factor $\sin\bigl(N+\frac12\bigr)u$ oscillates so fast that it averages everything out.

\begin{example}[the sawtooth]\label{ex:sawtooth}
Let $f(x)=\frac{\pi-x}2$ on $\left]0,2\pi\right[$, extended $2\pi$-periodically, with $f(0)=0$. Then $c_0(f)=0$ and, integrating by parts,
$c_n(f)=-\frac i{2n}$ for $n\ne0$, so $S_Nf(x)=\sum_{n=1}^N\frac{\sin nx}n$. The function is piecewise $C^1$, and $\frac{f(0^+)+f(0^-)}2=\frac{\frac\pi2-\frac\pi2}2=0$,
so Theorem~\ref{thm:dirichletfourier} gives, for every real $x$,
\[
\sum_{n\ge1}\frac{\sin nx}n=\frac{\pi-x}2\qquad(0<x<2\pi),
\]
and the sum is $0$ for $x\in2\pi\Z$.
This is the sum promised after Theorem~\ref{thm:dirichlet} in Chapter~\ref{ch:finite}. At $x=\frac\pi2$ it is Leibniz's series
$1-\frac13+\frac15-\cdots=\frac\pi4$ once more.
\end{example}

The two limits at $0$ are $\frac\pi2$ and $-\frac\pi2$. Their average is $0$, and that is the value of the series on the jump. The opening sums are this series at $x=1$ and at $x=2$.

\section{Fej\'er's theorem}

Pointwise convergence can fail for a continuous function. Exercise~\ref{ex:divergent} builds one. Fej\'er, in 1900, at the age of twenty, saw that the averages of the partial sums repair the failure.

\begin{definition}
The \emph{Fej\'er means} of $f$ are $\sigma_Nf=\frac1N\sum_{k=0}^{N-1}S_kf$ for $N\ge1$, trigonometric polynomials of degree at most $N-1$.
\end{definition}

For $f(x)=\cos x$ one has $S_0f=0$ and $S_kf=\cos x$ for $k\ge1$, so $\sigma_Nf=\frac{N-1}N\cos x$. At $N=2$ and $x=0$ the mean is $\frac12$, while $f(0)=1$. The average has moved halfway.

\begin{lemma}[Fej\'er kernel]\label{lem:fejerkernel}
$\sigma_Nf(x)=\frac1{2\pi}\int_{-\pi}^\pi f(x+u)F_N(u)\dd u$, where $F_N=\frac1N\sum_{k<N}D_k$ satisfies, for $u\notin2\pi\Z$,
\[
F_N(u)=\frac1N\Bigl(\frac{\sin(Nu/2)}{\sin(u/2)}\Bigr)^2.
\]
Hence $F_N\ge0$, $\frac1{2\pi}\int_{-\pi}^\pi F_N=1$, and $F_N(u)\le\frac1{N\sin^2(\delta/2)}$ for $\delta\le|u|\le\pi$.
\end{lemma}
\begin{proof}
The integral formula follows from Lemma~\ref{lem:dirichletkernel} by averaging. The identity $2\sin\bigl(k+\frac12\bigr)u\,\sin\frac u2=\cos ku-\cos(k+1)u$ makes the sum telescope:
\[
\sin\frac u2\sum_{k=0}^{N-1}\sin\Bigl(k+\frac12\Bigr)u=\frac{1-\cos Nu}2=\sin^2\frac{Nu}2.
\]
Divide by $N\sin^2\frac u2$. Each $D_k$ has mean $1$, so $F_N$ does too. For $\delta\le|u|\le\pi$ one has $\sin^2\frac u2\ge\sin^2\frac\delta2$.
\end{proof}

\begin{figure}[ht]
\centering
\fig{ch06-kernels}
\caption{The Dirichlet kernel $D_8$ (charcoal) and the Fej\'er kernel $F_8$ (coral) on $[-\pi,\pi]$. Both have mean $1$. The Dirichlet kernel
takes negative values and its oscillations decay slowly; the Fej\'er kernel is non-negative and concentrates at $0$.}
\label{fig:kernels}
\end{figure}

\begin{theorem}[Fej\'er, 1900]\label{thm:fejer}
If $f$ is continuous and $2\pi$-periodic, then $\sigma_Nf\to f$ uniformly on $\R$.
\end{theorem}
\begin{proof}
\emph{Setup.} The function $f$ is bounded, say by $M$, and uniformly continuous. It is continuous on $[-2\pi,2\pi]$, and periodic. Given $\varepsilon>0$, choose $\delta\in\left]0,\pi\right]$ so that $|f(x+u)-f(x)|\le\varepsilon$ for all $x$ and all $|u|\le\delta$.

\emph{Split.} Since $F_N\ge0$ has mean $1$, with $\Delta(u)=|f(x+u)-f(x)|$,
\[
\begin{aligned}|\sigma_Nf(x)-f(x)|&\le\frac1{2\pi}\int_{-\pi}^\pi\Delta(u)F_N(u)\dd u\eqbr\le\varepsilon+\frac{2M}{N\sin^2(\delta/2)}.\end{aligned}
\]
The integral over $|u|\le\delta$ contributes at most $\varepsilon$. The integral over $\delta\le|u|\le\pi$ contributes at most the second term, by the bound on $F_N$ in Lemma~\ref{lem:fejerkernel}.

\emph{Finish.} The bound does not depend on $x$. For $N$ large it is at most $2\varepsilon$.
\end{proof}

The positivity of $F_N$ is the whole difference with $D_N$. An average, with non-negative weights, of values close to $f(x)$ is close to $f(x)$. Three consequences follow at once.

\begin{corollary}\label{cor:fejer}
\begin{enumerate}[label=(\alph*)]
\item Trigonometric polynomials are uniformly dense among continuous $2\pi$-periodic functions.
\item A continuous $2\pi$-periodic function whose Fourier coefficients all vanish is zero.
\item \emph{Weierstrass's theorem}: every continuous function on $[a,b]$ is a uniform limit of polynomials.
\end{enumerate}
\end{corollary}
\begin{proof}
(a) is Theorem~\ref{thm:fejer}. (b) All the means $\sigma_Nf$ vanish, so their limit $f$ vanishes. (c) By an affine change of variable, take $[a,b]=[-1,1]$. The function $g(\theta)=f(\cos\theta)$ is continuous, even, and $2\pi$-periodic, so $c_{-n}(g)=c_n(g)$ and $\sigma_Ng(\theta)=\sum_na_n\cos n\theta$ for some coefficients $a_n$. From $\cos(n+1)\theta+\cos(n-1)\theta=2\cos\theta\cos n\theta$, induction gives $\cos n\theta=T_n(\cos\theta)$ for a polynomial $T_n$ of degree $n$. These are the Chebyshev polynomials. So
\[
\sup_{[-1,1]}\bigl|f-\sum_na_nT_n\bigr|=\sup_\theta|g-\sigma_Ng|\to0.
\]
\end{proof}

\section{Parseval's identity}

Bessel's inequality says the sum of $|c_n|^2$ cannot exceed the energy of $f$. Fej\'er's theorem upgrades the inequality to an equality, first for continuous $f$ and then for piecewise continuous $f$.

\begin{theorem}[Parseval]\label{thm:parseval}
For $f$ piecewise continuous, $\|f-S_Nf\|_2\to0$ and
\[
\frac1{2\pi}\int_{-\pi}^\pi|f(t)|^2\dd t=\sum_{n\in\Z}|c_n(f)|^2.
\]
More generally, $\langle f,g\rangle=\sum_nc_n(f)\overline{c_n(g)}$ for $f,g$ piecewise continuous.
\end{theorem}
\begin{proof}
\emph{Continuous case.} Let $f$ be continuous. The mean $\sigma_{N+1}f$ has degree at most $N$, so Proposition~\ref{prop:bestapprox} gives
\[
\begin{aligned}\|f-S_Nf\|_2&\le\|f-\sigma_{N+1}f\|_2\eqbr\le\sup|f-\sigma_{N+1}f|\to0.\end{aligned}
\]

\emph{Piecewise case.} Let $f$ be piecewise continuous, bounded by $M$, with $k$ points of discontinuity in a period. Given $\eta>0$ small, let $g$ equal $f$ except on intervals of length $\eta$ centred at these points, where $g$ is affine and joins the values of $f$ at the ends. Then $g$ is continuous, $|f-g|\le2M$, and $\|f-g\|_2^2\le\frac{4M^2k\eta}{2\pi}$, which is small for $\eta$ small. Since $S_Ng$ has degree at most $N$,
\[
\begin{aligned}\|f-S_Nf\|_2&\le\|f-S_Ng\|_2\eqbr\le\|f-g\|_2+\|g-S_Ng\|_2,\end{aligned}
\]
which is small for $N$ large, by the continuous case. The identity follows from $\|f-S_Nf\|_2^2=\|f\|_2^2-\sum_{|n|\le N}|c_n(f)|^2$.

\emph{Inner product.} For the general form, orthonormality gives $\langle S_Nf,S_Ng\rangle=\sum_{|n|\le N}c_n(f)\overline{c_n(g)}$. Cauchy--Schwarz gives
\[
\begin{aligned}\bigl|\langle f,g\rangle-\langle S_Nf,S_Ng\rangle\bigr|&\le\|f-S_Nf\|_2\|g\|_2\eqbr+\|S_Nf\|_2\|g-S_Ng\|_2,\end{aligned}
\]
which tends to $0$.
\end{proof}

Parseval's identity is Pythagoras with infinitely many orthogonal directions. The functions $e_n$ form a complete orthonormal system, and the squared length of $f$ is the sum of the squares of its coordinates. In physical language, the energy of a signal equals the energy of its spectrum. Exercise~\ref{ex:x2} spends the identity on two further proofs of $\zeta(2)=\frac{\pi^2}6$ and on $\zeta(4)=\frac{\pi^4}{90}$. Volume~8 places the identity in the theory of Hilbert spaces, where $L^2$ is the completion of the piecewise continuous functions for $\|\cdot\|_2$.

\section{The Gibbs phenomenon}

Let $f$ be the square wave: $f=1$ on $\left]0,\pi\right[$, $f=-1$ on $\left]-\pi,0\right[$, and $f(0)=f(\pi)=0$. Then $c_n(f)=\frac2{i\pi n}$ for $n$ odd and $c_n(f)=0$ for $n$ even, so
\[
S_{2N-1}f(x)=\frac4\pi\sum_{k=1}^N\frac{\sin(2k-1)x}{2k-1}.
\]
Figure~\ref{fig:rw06-square} draws the first three of these sums on $(0,\pi)$. Each one starts at $0$ and climbs toward the line of height $1$. Theorem~\ref{thm:dirichletfourier} says the sums tend to $f(x)$ at every $x$. The first arch of $S_1$ already reaches $\frac4\pi\approx1.27$, past the value $1$ of the wave.

\begin{figure}[ht]
\centering
\fig{ch06-rewrite-square}
\caption{$S_1$ (charcoal), $S_3$ (coral) and $S_5$ (rust) of the square wave on $(0,\pi)$, against the dashed line at height $1$.}
\label{fig:rw06-square}
\end{figure}

In 1898 the physicist Michelson built a machine that summed $80$ terms of such a series. He was surprised by bumps near the jumps that did not shrink. Gibbs explained them in 1899. Wilbraham had found them in 1848.

\begin{theorem}[Gibbs phenomenon]\label{thm:gibbs}
The first maximum of $S_{2N-1}f$ on $\left]0,\pi\right[$ is at $x_N=\frac\pi{2N}$, and
\[
S_{2N-1}f(x_N)\xrightarrow[N\to\infty]{}\frac2\pi\int_0^\pi\frac{\sin t}t\dd t=1.17897\ldots
\]
So the partial sums overshoot the jump of height $2$ by about $0.179$, nearly $9\%$ of the jump, however large $N$ is.
\end{theorem}
\begin{proof}
\emph{The derivative.} Differentiating and summing the cosines, the identity $2\sin x\sum_{k\le N}\cos(2k-1)x=\sin 2Nx$ gives
\[
\begin{aligned}\frac{\mathrm d}{\mathrm dx}S_{2N-1}f(x)&=\frac4\pi\sum_{k=1}^N\cos(2k-1)x\eqbr=\frac2\pi\cdot\frac{\sin2Nx}{\sin x}.\end{aligned}
\]
The derivative is positive on $\left]0,\frac\pi{2N}\right[$ and changes sign at $\frac\pi{2N}$. That point is the first maximum.

\emph{The Riemann sum.} Let $t_k=\frac{(2k-1)\pi}{2N}$. These are the midpoints of a subdivision of $[0,\pi]$ into $N$ intervals of length $\frac\pi N$, and
\[
S_{2N-1}f(x_N)=\frac4\pi\sum_{k=1}^N\frac{\sin t_k}{2k-1}=\frac2\pi\cdot\frac\pi N\sum_{k=1}^N\frac{\sin t_k}{t_k}.
\]
The sum is a Riemann sum for $\frac{\sin t}t$ on $[0,\pi]$, with the value $1$ at $t=0$. It tends to $\frac2\pi\operatorname{Si}(\pi)$. Figure~\ref{fig:si} records $\operatorname{Si}(\pi)\approx1.85194$.
\end{proof}

\begin{figure}[ht]
\centering
\fig{ch06-gibbs}
\caption{The square wave and its partial sums $S_5$ (ochre) and $S_{39}$ (charcoal). The overshoot near the jump moves closer to it but
keeps its height, about $1.179$ (dashed).}
\label{fig:gibbs}
\end{figure}

\begin{think}
Theorem~\ref{thm:dirichletfourier} says $S_Nf(x)\to f(x)$ at every point, and Theorem~\ref{thm:gibbs} says $\sup_{]0,\pi[}S_Nf$ does not tend to $1$. Why is there no
contradiction? What does the Fej\'er mean $\sigma_Nf$ do near the jump?
\end{think}

There is no contradiction, because the bump moves. For each fixed $x$ in $(0,\pi)$ the sums eventually sit near $1$. The place where they reach about $1.179$ slides toward the jump, so the supremum on the whole interval stays large. The Fej\'er means do not overshoot. The kernel $F_N$ is non-negative and has mean $1$, so $\sigma_Nf$ stays between the bounds of $f$.

The proportion is the same at every jump, and the comparison uses the sawtooth, which has one jump in each period. The sawtooth $s$ of Example~\ref{ex:sawtooth} jumps by $\pi$ at $0$: the right-hand limit is $\frac\pi2$ and the left-hand limit is $-\frac\pi2$. Its partial sum at $\frac\pi N$ is
\[
\sum_{n=1}^N\frac{\sin(n\pi/N)}n=\frac\pi N\sum_{n=1}^N\frac{\sin t_n}{t_n},\qquad t_n=\frac{n\pi}N.
\]
This is again a Riemann sum, and it tends to $\operatorname{Si}(\pi)$. The value just to the right of the jump is $\frac\pi2$, so the sawtooth overshoots by $\operatorname{Si}(\pi)-\frac\pi2\approx0.281$. Divided by the jump $\pi$, that overshoot is $\frac1\pi\operatorname{Si}(\pi)-\frac12\approx0.0895$. The square wave, on a jump of height $2$, overshoots by the same fraction of its jump.

Now let $f$ be piecewise $C^1$, with a jump of size $J$ at $a$. Subtract $\frac J\pi\, s(x-a)$. The multiple of the sawtooth has one jump, of size $J$, because $s$ jumps by $\pi$. Do the same at the other jumps. What remains is continuous and piecewise $C^1$, so its Fourier series converges uniformly, by Exercise~\ref{ex:abs}. A uniform limit has no room for a bump that stays a fixed height above the function. Every jump therefore produces an overshoot of $\bigl(\frac1\pi\operatorname{Si}(\pi)-\frac12\bigr)J\approx0.0895\,J$. Signal processing uses a smooth window, because the Fej\'er means stay inside the bounds of $f$.

\section{Equidistribution}

A real sequence $(u_n)$ is \emph{equidistributed modulo $1$} when, for all $0\le a<b\le1$, the proportion of $n\le N$ with fractional part $\{u_n\}\in[a,b[$ tends to $b-a$. Fourier analysis on the circle $\R/\Z$ turns the definition into one exponential sum for each frequency.

\begin{theorem}[Weyl's criterion, 1916]\label{thm:weyl}
$(u_n)$ is equidistributed modulo $1$ if and only if, for every integer $h\ne0$,
\[
\frac1N\sum_{n=1}^Ne^{2\pi ihu_n}\xrightarrow[N\to\infty]{}0.
\]
\end{theorem}
\begin{proof}
\emph{Good functions.} For a $1$-periodic function $F$ let $A_N(F)=\frac1N\sum_{n\le N}F(u_n)$, and say that $F$ is \emph{good} when $A_N(F)\to\int_0^1 F$. The good functions form a vector space. Equidistribution says that the $1$-periodic extensions of the indicators $\mathbf 1_{[a,b[}$ are good. The criterion says that $e^{2\pi ihx}$, for $h\ne0$, is good, because its integral is $0$. Constants are always good.

\emph{From indicators to exponentials.} Suppose every $\mathbf 1_{[a,b[}$ is good. Then every step function is good. A continuous $1$-periodic $F$ is uniformly continuous, so for $\varepsilon>0$ there is a step function $s$ with $|F-s|\le\varepsilon$ everywhere. Then $\bigl|A_N(F)-\int F\bigr|\le2\varepsilon+\bigl|A_N(s)-\int s\bigr|$, and $F$ is good. In particular the real and imaginary parts of $e^{2\pi ihx}$ are good, which is the criterion.

\emph{From exponentials back.} Suppose the criterion holds. Then every trigonometric polynomial in the functions $e^{2\pi ihx}$ is good. By Fej\'er's theorem, applied to $\theta\mapsto F\bigl(\frac\theta{2\pi}\bigr)$, every continuous $1$-periodic $F$ is a uniform limit of such polynomials. The same $2\varepsilon$ argument shows that $F$ is good. Now fix $[a,b[$ and $\varepsilon>0$. There are continuous $1$-periodic functions $F_-\le\mathbf 1_{[a,b[}\le F_+$ with $\int F_+-\int F_-\le4\varepsilon$. Take $F_+$ equal to $1$ on $[a,b]$, equal to $0$ at distance more than $\varepsilon$ from that interval, and affine in between. Build $F_-$ the same way inside $[a,b]$, or take $F_-=0$ if $b-a<2\varepsilon$. Since $A_N$ is monotone,
\[
\begin{aligned}\int F_-&\le\liminf A_N(\mathbf 1_{[a,b[})\eqbr\le\limsup A_N(\mathbf 1_{[a,b[})\le\int F_+,\end{aligned}
\]
and both bounds lie within $4\varepsilon$ of $b-a$.
\end{proof}

\begin{corollary}[Bohl, Sierpi\'nski, Weyl, 1909--1910]\label{cor:nalpha}
If $\alpha$ is irrational, $(n\alpha)$ is equidistributed modulo $1$.
\end{corollary}
\begin{proof}
For $h\ne0$ one has $h\alpha\notin\Z$, so $q=e^{2\pi ih\alpha}\ne1$ and
\[
\Bigl|\sum_{n=1}^Nq^n\Bigr|=\Bigl|\frac{q^{N+1}-q}{q-1}\Bigr|\le\frac2{|q-1|}.
\]
The bound does not depend on $N$, so the average tends to $0$. Weyl's criterion applies.
\end{proof}

\begin{figure}[ht]
\centering
\fig{ch06-weyl}
\caption{The fractional parts $\{n\varphi\}$ for $n=1,\dots,60$, with $\varphi$ the golden ratio, plotted at height $n$. Each new point
falls into one of the largest gaps left by the previous ones. In this picture the gaps take at most three lengths.}
\label{fig:weyl}
\end{figure}

For rational $\alpha=\frac pq$ the fractional parts take only $q$ values, so irrationality is necessary. Figure~\ref{fig:rw06-equi} puts the two behaviours side by side for $n=1,\ldots,20$. Take $\alpha=\frac{\sqrt5-1}2$, which is irrational, because $\sqrt5$ is. The twenty fractional parts $\{n\alpha\}$ are spread along $[0,1]$, and exactly ten of them lie in $[0,\frac12)$. Take instead $\frac{n^2}5$. The square of an integer is $0$, $1$, or $4$ modulo $5$, so $\{n^2/5\}$ lands only on $0$, $\frac15$, and $\frac45$, with $4$, $8$, and $8$ hits. A short cycle is order. An irrational rotation fills the interval the way a uniform random sample would. Exercise~\ref{ex:benford} spends the corollary on the first digits of the powers of $2$. Volume~14 estimates the sums $\sum_{n\le N}e^{2\pi i\alpha n^k}$, which decide whether $(n^k\alpha)$ is equidistributed, on its way to Waring's problem. That is a preview. Nothing here depends on it.

\begin{figure}[ht]
\centering
\fig{ch06-rewrite-equi}
\caption{For $n=1,\ldots,20$, the parts $\{n\alpha\}$ with $\alpha=(\sqrt5-1)/2$ are spread on $[0,1]$ (charcoal), while $\{n^2/5\}$ sits on $0$, $\frac15$ and $\frac45$ with $4$, $8$ and $8$ hits (rust).}
\label{fig:rw06-equi}
\end{figure}

\section{Several variables and inversion on \texorpdfstring{$\Z^d$}{Zd}}

For $k\in\Z^d$ and $t\in\R^d$ write $k\cdot t=k_1t_1+\cdots+k_dt_d$. Integrals of continuous functions over the cube $Q=[-\pi,\pi]^d$ are iterated integrals. By Theorem~\ref{thm:fubinirect} and induction on $d$, the order does not matter. A uniform limit of continuous functions on $Q$ can be integrated term by term, one variable at a time.

\begin{proposition}[Fourier inversion on $\Z^d$]\label{prop:inversionZd}
Let $(a_k)_{k\in\Z^d}$ be complex numbers with $\sum_k|a_k|<\infty$, and $\varphi(t)=\sum_ka_ke^{ik\cdot t}$. The series converges uniformly on $\R^d$,
$\varphi$ is continuous, and for every $k\in\Z^d$
\[
a_k=\frac1{(2\pi)^d}\int_Q\varphi(t)e^{-ik\cdot t}\dd t.
\]
\end{proposition}
\begin{proof}
Enumerate $\Z^d$. The terms are bounded by $|a_k|$, so the series converges uniformly, by the Weierstrass test, Lemma~\ref{lem:uniform}. It converges in any order to the same sum, by Lemma~\ref{lem:packets}, and $\varphi$ is continuous. For $k,l\in\Z^d$ the iterated integral factorises:
\[
\int_Qe^{i(l-k)\cdot t}\dd t=\prod_{j=1}^d\int_{-\pi}^\pi e^{i(l_j-k_j)t_j}\dd t_j,
\]
which is $(2\pi)^d$ if $l=k$ and $0$ otherwise. Multiply the series by $e^{-ik\cdot t}$. The product still converges uniformly. Integrate term by term.
\end{proof}

\begin{example}[return probabilities of random walks]
Let $S_n=X_1+\cdots+X_n$ be a random walk on $\Z^d$ with independent steps distributed like $X$, and let $\varphi(t)=E\bigl[e^{it\cdot X}\bigr]=\sum_kP(X=k)e^{ik\cdot t}$
be its characteristic function. By independence of the steps (Volume~4, Chapter~2), $E\bigl[e^{it\cdot S_n}\bigr]=\varphi(t)^n=\sum_kP(S_n=k)e^{ik\cdot t}$, and the
probabilities are summable, so Proposition~\ref{prop:inversionZd} gives
\[
P(S_n=0)=\frac1{(2\pi)^d}\int_Q\varphi(t)^n\dd t.
\]
This is the formula used in Volume~4, Chapter~8, to decide when a walk returns to its start infinitely often.
\end{example}

The closed form of the characteristic function is order. Whether the integral of its powers tends to zero fast enough for the walk to escape is the other half of the question, and Volume~4 answers it.

\section{The isoperimetric inequality}

Among closed curves of a given length, which one encloses the largest area? The Greeks knew the answer, the circle. A proof waited until the nineteenth century. Hurwitz found the shortest one in 1901, and the tool is Parseval's identity.

Let $z(s)=x(s)+iy(s)$, for $s\in[0,2\pi]$, be a closed curve of length $L$, of class $C^1$, parametrised proportionally to arc length, so that $|z'(s)|=\frac L{2\pi}$ for all $s$. If the curve is simple and traversed counterclockwise, the area it encloses is
\[
\begin{aligned}A&=\frac12\int_0^{2\pi}\bigl(xy'-yx'\bigr)\dd s\eqbr=\frac12\operatorname{Im}\int_0^{2\pi}\overline{z(s)}z'(s)\dd s.\end{aligned}
\]
This is Green's formula. For a polygon it is the shoelace formula. Here we take it as the definition of the area.

\begin{theorem}[isoperimetric inequality]\label{thm:isoperimetric}
$4\pi A\le L^2$, with equality if and only if the curve is a circle.
\end{theorem}
\begin{proof}
\emph{Coefficients.} Extend $z$ to a $2\pi$-periodic function. It is continuous, and $z'$ is piecewise continuous. Let $c_n=c_n(z)$. Integration by parts, using $z(-\pi)=z(\pi)$, gives $c_n(z')=inc_n$.

\emph{Length and area.} Parseval for $z'$ gives
\[
\Bigl(\frac L{2\pi}\Bigr)^2=\frac1{2\pi}\int_{-\pi}^\pi|z'|^2=\sum_nn^2|c_n|^2.
\]
The general form of Parseval gives $\frac1{2\pi}\int z'\bar z=\sum_n(inc_n)\overline{c_n}=i\sum_nn|c_n|^2$, so $A=\pi\sum_nn|c_n|^2$. Hence
\[
L^2-4\pi A=4\pi^2\sum_{n\in\Z}\bigl(n^2-n\bigr)|c_n|^2\ge0,
\]
since $n^2\ge n$ for every integer $n$, with equality only for $n\in\{0,1\}$.

\emph{Equality.} Equality holds if and only if $c_n=0$ for $n\ne0,1$, that is, $z(s)=c_0+c_1e^{is}$. That is a circle.
\end{proof}

The coefficients are the shape. Every mode except $n=0$ and $n=1$ adds to the length faster than it adds to the area, and the inequality is the sum of those non-negative contributions.

\section{What we proved, and what comes next}

Fourier's claim holds in three senses. At a point, the series of a piecewise $C^1$ function converges to the function, and to the middle of a jump (Theorem~\ref{thm:dirichletfourier}). Uniformly, the averages of the partial sums of a continuous function converge to it, though the partial sums may diverge (Theorem~\ref{thm:fejer}, Exercise~\ref{ex:divergent}). In mean square, every piecewise continuous function has a convergent series, and the energy of the function is the sum of the energies of its harmonics (Theorem~\ref{thm:parseval}). A sum of smooth waves can jump. At the jump the series takes the middle value, and near the jump the partial sums overshoot by about $9\%$ of the jump, for every $N$ (Theorem~\ref{thm:gibbs}).

Three applications used nothing beyond these theorems. Irrational rotations are equidistributed (Corollary~\ref{cor:nalpha}). The probability that a random walk is back at the start is an integral (Proposition~\ref{prop:inversionZd}). The circle maximises area for a given length (Theorem~\ref{thm:isoperimetric}).

The next two chapters keep the orthogonality and change the group. Chapter~\ref{ch:fft} replaces the circle by $N$ equally spaced points, where every sum is finite and the work can be done in $N\log N$ steps. Chapter~\ref{ch:transform} replaces the circle by the line, where the series becomes an integral, and the Poisson summation formula ties the line back to the circle.

\begin{beauty}
One identity, Pythagoras for the functions $e^{inx}$, gives the energy of a signal, two values of $\zeta$, the circle as the solution of the isoperimetric problem, and the probability that a random walk comes home.
\end{beauty}

\section{Exercises}

\begin{exercise}[$\star$]\label{ex:x2}
Compute the Fourier coefficients of the $2\pi$-periodic functions equal to $x$ and to $x^2$ on $\left]-\pi,\pi\right]$. Deduce $\zeta(2)=\frac{\pi^2}6$ twice,
from Dirichlet's theorem at $x=\pi$ for $x^2$ and from Parseval for $x$, and $\zeta(4)=\frac{\pi^4}{90}$ from Parseval for $x^2$.
\end{exercise}

\begin{exercise}[$\star\star$]\label{ex:abs}
Let $f$ be continuous, $2\pi$-periodic and piecewise $C^1$. Show that $c_n(f')=inc_n(f)$, that $\sum_n|c_n(f)|<\infty$, and that $S_Nf\to f$
uniformly.
\end{exercise}

\begin{exercise}[$\star\star$]\label{ex:heat}
Let $f$ be as in Exercise~\ref{ex:abs}, and for $t>0$ let $u(t,x)=\sum_nc_n(f)e^{-n^2t}e^{inx}$. Show that $\partial_tu=\partial_x^2u$ for $t>0$, and that
$u(t,\cdot)\to f$ uniformly as $t\to0^+$. Which harmonics of the initial temperature survive longest?
\end{exercise}

\begin{exercise}[$\star\star$]\label{ex:wirtinger}
(Wirtinger's inequality.) Let $f$ be real, $2\pi$-periodic and of class $C^1$, with $\int_{-\pi}^\pi f=0$. Show that $\int_{-\pi}^\pi f^2\le\int_{-\pi}^\pi f'^2$, and find
the cases of equality.
\end{exercise}

\begin{exercise}[$\star\star$]\label{ex:benford}
Show that the proportion of $n\le N$ for which $2^n$ starts with the digit $d$ tends to $\log_{10}\bigl(1+\frac1d\bigr)$. How often does $2^n$ start
with $1$? With $7$?
\end{exercise}

\begin{exercise}[$\star\star\star$]\label{ex:divergent}
(A continuous function whose Fourier series diverges at $0$, after Fej\'er.)
\begin{enumerate}[label=(\alph*)]
\item Show that $\bigl|\sum_{j=1}^n\frac{\sin jx}j\bigr|\le1+\pi$ for all $n$ and $x$.
\item Let $P_n(x)=2\sin(2nx)\sum_{j=1}^n\frac{\sin jx}j$. Show that $|P_n|\le2+2\pi$, that $P_n(x)=\sum_{j=1}^n\frac{\cos(2n-j)x-\cos(2n+j)x}j$, and that $P_n(0)=0$.
\item Let $n_k=2^{k^3}$ and $f=\sum_{k\ge1}\frac1{k^2}P_{n_k}$. Show that $f$ is continuous and that $S_{2n_k-1}f(0)\ge k\ln2$.
\end{enumerate}
\end{exercise}

\begin{exercise}[$\star\star\star$]\label{pb:flat}
For signs $\varepsilon_0,\dots,\varepsilon_n\in\{-1,1\}$, let $P(z)=\sum_{k=0}^n\varepsilon_kz^k$ and
$\|P\|=\max_{|z|=1}|P(z)|$.
\begin{enumerate}[label=(\alph*)]
\item Show that $\|P\|\ge\sqrt{n+1}$ for every choice of signs.
\item Show that for every $n\ge2$ some choice of signs gives
\[
\|P\|\le2\sqrt{(n+1)\log\bigl(5n(n+1)\bigr)}+\frac\pi2.
\]
\end{enumerate}
\end{exercise}

\begin{exercise}[$\star\star$]\label{pb:rudin}
(Golay 1951, Shapiro 1952, Rudin 1959.) Let $P_0=Q_0=1$ and, for $k\ge0$,
\[
P_{k+1}(z)=P_k(z)+z^{2^k}Q_k(z),\qquad Q_{k+1}(z)=P_k(z)-z^{2^k}Q_k(z).
\]
\begin{enumerate}[label=(\alph*)]
\item Show that $P_k$ and $Q_k$ have degree $2^k-1$ and all their coefficients in $\{-1,1\}$.
\item Show that $|P_k(z)|^2+|Q_k(z)|^2=2^{k+1}$ for $|z|=1$, and deduce that $\|P_k\|\le\sqrt2\sqrt{2^k}$: within a factor
$\sqrt2$ of the lower bound of Problem~\ref{pb:flat}(a) (Figure~\ref{fig:rudin}).
\item Show that the coefficient of $z^m$ in $P_k$ (for $m<2^k$) is $(-1)^{e(m)}$, where $e(m)$ is the number of pairs of
adjacent $1$s in the binary expansion of $m$.
\end{enumerate}
\end{exercise}

\begin{exercise}[$\star\star\star$]\label{pb:siegel}
\emph{Minkowski's theorem by Fourier analysis (Siegel, 1935).} Let $K\subset\R^n$ be a bounded convex set with $K=-K$, and suppose
that $K$ contains no point of $\Z^n$ other than $0$. This exercise integrates over $K$, which the volume has not defined for a
general convex set. Admitted here: $K$ has a volume, and the integral of a continuous function over $\frac12K$ and its
translates is additive and invariant under translation, as iterated integrals over boxes are. Volume~7, Theorem~7.11, builds Lebesgue measure on $\R^n$, and a bounded convex set is a Borel set, so it has that volume.
\begin{enumerate}[label=(\alph*)]
\item Show that the translates $\frac12K+m$, $m\in\Z^n$, are pairwise disjoint.
\item Let $\varphi(x)=\sum_{m\in\Z^n}1_{\frac12K}(x+m)$, a $\Z^n$-periodic function. Show that $\varphi^2=\varphi$, and that its Fourier coefficients
on the torus $[0,1[^n$ are $\hat\varphi(0)=2^{-n}\operatorname{vol}K$ and $\hat\varphi(k)=\int_{\frac12K}e(-k\cdot x)\dd x$, where $e(t)=e^{2\pi it}$.
\item Using Bessel's inequality on the torus, deduce that
\[
2^{-n}\operatorname{vol}K\bigl(1-2^{-n}\operatorname{vol}K\bigr)\ge\sum_{k\ne0}\Bigl|\int_{\frac12K}e(-k\cdot x)\dd x\Bigr|^2,
\]
and in particular Minkowski's theorem: a symmetric convex body of volume greater than $2^n$ contains a non-zero integer point.
\item Check that for the open cube $K=\left]-1,1\right[^n$ both sides vanish.
\end{enumerate}
\end{exercise}

\begin{exercise}[$\star$]\label{ex:rw06a}
The square wave equals $1$ on $(0,\pi)$, and
\[
S_{2N-1}f(x)=\frac4\pi\sum_{k=1}^N\frac{\sin(2k-1)x}{2k-1}.
\]
Evaluate $S_1f$, $S_3f$ and $S_5f$ at $x=\frac\pi2$. The series converges to $1$ there. Which of the three values lie above $1$?
\end{exercise}

\begin{exercise}[$\star\star$]\label{ex:rw06b}
The sawtooth of Example~\ref{ex:sawtooth} has $c_0=0$ and $c_n=-\frac i{2n}$ for $n\ne0$.
\begin{enumerate}[label=(\alph*)]
\item Check that $\frac1{2\pi}\int_0^{2\pi}\bigl(\frac{\pi-x}2\bigr)^2\dd x=\frac{\pi^2}{12}$.
\item Parseval says this equals $\sum_{n\ne0}\frac1{4n^2}$. Deduce $\sum_{n\ge1}\frac1{2n^2}=\frac{\pi^2}{12}$ from $\zeta(2)=\frac{\pi^2}6$.
\item Compute $\sum_{n=1}^5\frac1{2n^2}$ as a fraction, and the gap between that fraction and $\frac{\pi^2}{12}$.
\end{enumerate}
\end{exercise}

\begin{exercise}[$\star\star$]\label{ex:rw06c}
Let $\alpha=\frac{\sqrt5-1}2$.
\begin{enumerate}[label=(\alph*)]
\item Show that $\alpha$ is irrational.
\item Among $n=1,\ldots,20$, how many fractional parts $\{n\alpha\}$ lie in $\bigl[0,\frac12\bigr)$?
\item Among $n=1,\ldots,20$, how many fractional parts $\{n^2/5\}$ lie in $\bigl[0,\frac12\bigr)$? Show that, as the upper limit grows, this proportion tends to $\frac35$.
\end{enumerate}
\end{exercise}

\begin{exercise}[$\star\star\star$]\label{ex:rw06d}
\begin{enumerate}[label=(\alph*)]
\item Show that $(n\sqrt2)$ is equidistributed modulo $1$.
\item Let $q=e^{2\pi i\sqrt2}$. Check that $\bigl|\sum_{n=1}^{10}q^n\bigr|\le\frac2{|q-1|}$, and give both sides to three decimals.
\item Show that $(n/3)$ is not equidistributed modulo $1$. Among $n=1,\ldots,30$, how many fractional parts $\{n\sqrt2\}$ lie in $[0,0.1)$, and how many fractional parts $\{n/3\}$ lie in $\bigl[0,\frac12\bigr)$?
\end{enumerate}
\end{exercise}

\begin{figure}[ht]
\centering
\fig{ch01-rudin}
\caption{$|P(e^{i\theta})|$ for the Rudin--Shapiro polynomial $P_6$ (charcoal, $64$ coefficients) and for $64$ random signs (rust).
The first never exceeds $\sqrt{128}\approx11.31$; the second reaches $15.3$. Both have mean square $64$.}
\label{fig:rudin}
\end{figure}

\section{Hints}

\begin{hint}{ex:x2}
Integrate by parts; $c_n(x)=\frac{i(-1)^n}n$ and $c_n(x^2)=\frac{2(-1)^n}{n^2}$ for $n\ne0$.
\end{hint}
\begin{hint}{ex:abs}
Integrate by parts on each piece; the boundary terms cancel by continuity and periodicity. Then Cauchy--Schwarz:
$\sum_{n\ne0}|c_n|=\sum_{n\ne0}\frac1{|n|}\cdot|nc_n|$.
\end{hint}
\begin{hint}{ex:heat}
For $t\ge t_0>0$, the derivatives of the terms are bounded by $n^2|c_n|e^{-n^2t_0}$, which is summable. For the limit, split the series
into $|n|\le M$ and $|n|>M$.
\end{hint}
\begin{hint}{ex:wirtinger}
Parseval for $f$ and $f'$, with $c_0(f)=0$.
\end{hint}
\begin{hint}{ex:benford}
$2^n$ starts with $d$ when $\{n\log_{10}2\}\in\bigl[\log_{10}d,\log_{10}(d+1)\bigr[$. Is $\log_{10}2$ rational?
\end{hint}
\begin{hint}{ex:divergent}
(a) For $0<x<\pi$, treat $j\le\frac1x$ with $|\sin jx|\le jx$, and $j>\frac1x$ with summation by parts (Proposition~\ref{prop:abel}) and
$\bigl|\sum_{j=a}^b\sin jx\bigr|\le\frac1{\sin(x/2)}\le\frac\pi x$. (c) Which frequencies does each $P_{n_k}$ contain?
\end{hint}

\begin{hint}{pb:flat}
(a) Parseval. (b) Random signs, Hoeffding's inequality for the real and imaginary parts at $M=n(n+1)$ equally spaced points,
and $|P'(z)|\le\frac{n(n+1)}2$ between them.
\end{hint}

\begin{hint}{pb:rudin}
The parallelogram law $|u+v|^2+|u-v|^2=2|u|^2+2|v|^2$. For (c), find and prove a formula for the coefficients of $Q_k$ too.
\end{hint}

\begin{hint}{pb:siegel}
(a) If $x+m=y+m'$ with $x,y\in\frac12K$, then $m'-m\in K$. (b) Unfold the integral over the torus into an integral over $\R^n$. (c) Bessel:
$\int|\varphi|^2\ge\sum_{k\in F}|\hat\varphi(k)|^2$ for every finite $F$.
\end{hint}

\begin{hint}{ex:rw06a}
At $x=\frac\pi2$ one has $\sin\bigl((2k-1)\frac\pi2\bigr)=(-1)^{k-1}$. The three sums are $1$, $1-\frac13$, and $1-\frac13+\frac15$, each times $\frac4\pi$.
\end{hint}
\begin{hint}{ex:rw06b}
The substitution $u=\pi-x$ turns the integral into $\frac1{2\pi}\int_{-\pi}^\pi\frac{u^2}4\dd u$. For (c), the common denominator of $1,4,9,16,25$ is $3600$.
\end{hint}
\begin{hint}{ex:rw06c}
If $\alpha$ were rational then $\sqrt5=2\alpha+1$ would be rational. The squares modulo $5$ are $0$, $1$ and $4$.
\end{hint}
\begin{hint}{ex:rw06d}
$\sqrt2$ is irrational, so Corollary~\ref{cor:nalpha} applies. The sum of the geometric progression is $\frac{q(q^{10}-1)}{q-1}$. The parts $\{n/3\}$ cycle through $\frac13$, $\frac23$, $0$.
\end{hint}

\section{Solutions}

\begin{solution}{ex:x2}
\idea Two integrations by parts give the coefficients. Then either evaluate where every cosine has sign $(-1)^n$, or square and integrate.
\writeup For $n\ne0$,
\[
\int_{-\pi}^\pi xe^{-inx}\dd x=\Bigl[\frac{xe^{-inx}}{-in}\Bigr]_{-\pi}^\pi=\frac{2\pi(-1)^n}{-in},
\]
so $c_n(x)=\frac{i(-1)^n}n$, and $c_0(x)=0$. The boundary term in $\int x^2e^{-inx}\dd x$ vanishes, and what remains is $\frac2{in}\int xe^{-inx}\dd x$. Hence $c_n(x^2)=\frac2{in}\cdot\frac{i(-1)^n}n=\frac{2(-1)^n}{n^2}$, and $c_0(x^2)=\frac{\pi^2}3$. In real form,
\[
x^2=\frac{\pi^2}3+4\sum_{n\ge1}\frac{(-1)^n}{n^2}\cos nx\qquad(|x|\le\pi),
\]
with equality everywhere by Theorem~\ref{thm:dirichletfourier}, since the periodic extension of $x^2$ is continuous and piecewise $C^1$. At $x=\pi$, one has $\cos n\pi=(-1)^n$ and $\pi^2=\frac{\pi^2}3+4\zeta(2)$, so $\zeta(2)=\frac{\pi^2}6$. Parseval for $x$ reads $\frac1{2\pi}\int x^2=\frac{\pi^2}3=\sum_{n\ne0}\frac1{n^2}=2\zeta(2)$, the same value again. Parseval for $x^2$ reads $\frac1{2\pi}\int x^4=\frac{\pi^4}5=\frac{\pi^4}9+8\zeta(4)$, so
\[
\zeta(4)=\frac{\pi^4}8\Bigl(\frac15-\frac19\Bigr)=\frac{\pi^4}{90}.
\]
\end{solution}

\begin{solution}{ex:abs}
\idea Differentiation multiplies the $n$th coefficient by $in$. The coefficients of $f'$ are square-summable, so those of $f$ decay fast enough to be summable.
\writeup Let $-\pi=a_0<\cdots<a_m=\pi$ be the pieces. On each piece,
\[
\int_{a_j}^{a_{j+1}}f'e^{-inx}\dd x=\bigl[fe^{-inx}\bigr]_{a_j}^{a_{j+1}}+in\int_{a_j}^{a_{j+1}}fe^{-inx}\dd x.
\]
Since $f$ is continuous, the boundary terms telescope to $f(\pi)e^{-in\pi}-f(-\pi)e^{in\pi}=0$, by periodicity. So $c_n(f')=inc_n(f)$. Cauchy--Schwarz in $\ell^2$, and Bessel for $f'$, give
\[
\begin{aligned}\sum_{n\ne0}|c_n(f)|&\le\Bigl(\sum_{n\ne0}\frac1{n^2}\Bigr)^{1/2}\Bigl(\sum_n|c_n(f')|^2\Bigr)^{1/2}\eqbr\le\sqrt{2\zeta(2)}\,\|f'\|_2.\end{aligned}
\]
The Weierstrass test makes $S_Nf$ converge uniformly, to a continuous function $g$. Integrating term by term gives $c_n(g)=c_n(f)$, so $g=f$ by Corollary~\ref{cor:fejer}(b).
\end{solution}

\begin{solution}{ex:heat}
\idea Each harmonic $e^{inx}$ solves the heat equation once it is multiplied by $e^{-n^2t}$. The same factor makes every series converge quickly for $t>0$.
\writeup Each term $c_ne^{-n^2t}e^{inx}$ satisfies $\partial_t=\partial_x^2$, since both sides multiply the term by $-n^2$. On $t\ge t_0$, the term and its derivatives in $t$ and in $x$, up to order two, are bounded by $(1+n^2)|c_n|e^{-n^2t_0}$, which is summable. Lemma~\ref{lem:uniform}(d), applied in each variable, lets us differentiate term by term, and $u$ solves the equation. For the limit, Exercise~\ref{ex:abs} gives $|u(t,x)-f(x)|\le\sum_n|c_n|(1-e^{-n^2t})$. Given $\varepsilon$, choose $M$ with $\sum_{|n|>M}|c_n|\le\varepsilon$. The sum is then at most $(1-e^{-M^2t})\sum|c_n|+\varepsilon$, which is at most $2\varepsilon$ for $t$ small. The harmonic $n$ decays like $e^{-n^2t}$. High frequencies die first. The constant $c_0$, the mean temperature, survives. After a while only $c_0+c_{\pm1}e^{-t}e^{\pm ix}$ is visible.
\end{solution}

\begin{solution}{ex:wirtinger}
\idea On the Fourier side, differentiation multiplies the $n$th coefficient by $in$, and $|n|\ge1$ once $c_0$ is gone.
\writeup By Exercise~\ref{ex:abs}, $c_n(f')=inc_n(f)$, and $c_0(f)=0$. Parseval gives
\[
\frac1{2\pi}\int f'^2=\sum n^2|c_n|^2\ge\sum|c_n|^2=\frac1{2\pi}\int f^2.
\]
Equality holds if and only if $c_n=0$ for $|n|\ge2$, that is, $f(x)=a\cos x+b\sin x$.
\end{solution}

\begin{solution}{ex:benford}
\idea The first digit of $2^n$ is read from the fractional part of $n\log_{10}2$, and those fractional parts are equidistributed.
\writeup The power $2^n$ starts with the digit $d$ if and only if $d\cdot10^m\le2^n<(d+1)10^m$ for some integer $m$, if and only if $\{n\log_{10}2\}\in[\log_{10}d,\log_{10}(d+1)[$. If $\log_{10}2=\frac pq$ with $q\ge1$, then $2^q=10^p$. That is impossible, because $5$ does not divide $2^q$, and $p=0$ would give $2^q=1$. Corollary~\ref{cor:nalpha} says the proportion tends to $\log_{10}(d+1)-\log_{10}d=\log_{10}\bigl(1+\frac1d\bigr)$. About $30.1\%$ of the powers of $2$ start with $1$, and about $5.8\%$ start with $7$. This is Benford's law. It shows up in tables of data that span several orders of magnitude.
\end{solution}

\begin{solution}{ex:divergent}
\idea Build blocks of trigonometric polynomials that stay bounded, while their partial sums at $0$ grow like $\ln n$. Place the blocks so far apart in frequency that they do not interact.
\writeup \emph{The sine bound.} The sum in (a) is odd and $2\pi$-periodic in $x$, so take $0<x<\pi$, and let $m=\lfloor\frac1x\rfloor$. For $j\le m$, one has $\bigl|\frac{\sin jx}j\bigr|\le x$, so these terms contribute at most $mx\le1$. For $j>m$,
\[
\sum_{j=m+1}^b\sin jx=\frac{\cos\bigl(m+\frac12\bigr)x-\cos\bigl(b+\frac12\bigr)x}{2\sin(x/2)}
\]
has absolute value at most $\frac1{\sin(x/2)}\le\frac\pi x$, because $\sin\frac x2\ge\frac x\pi$ on $[0,\pi]$. Abel summation with the decreasing weights $\frac1j$, as in the proof of Theorem~\ref{thm:dirichlet}, bounds the rest by $\frac1{m+1}\cdot\frac\pi x\le\pi$, since $m+1>\frac1x$.

\emph{The block.} Part (b) follows from (a). The product $2\sin a\sin b=\cos(a-b)-\cos(a+b)$ gives the expansion. Its frequencies lie in $[n,3n]$, and $P_n(0)=0$ because $\sin0=0$.

\emph{The partial sum.} The series in (c) converges uniformly, by (b) and $\sum\frac1{k^2}<\infty$, so $f$ is continuous and its coefficients are obtained term by term. Since $3n_k<n_{k+1}$, the blocks occupy disjoint ranges of frequencies. The partial sum $S_{2n_k-1}$ keeps every block $i<k$, whose frequencies are at most $3n_i<n_k$. It keeps none of the blocks $i>k$, whose frequencies are at least $n_{k+1}>2n_k$. From block $k$ it keeps exactly the terms $\frac{\cos(2n_k-j)x}j$. At $x=0$ the blocks $i<k$ contribute $\frac{P_{n_i}(0)}{i^2}=0$, so
\[
S_{2n_k-1}f(0)=\frac1{k^2}\sum_{j=1}^{n_k}\frac1j\ge\frac{\ln n_k}{k^2}=k\ln2\to\infty.
\]
The Fourier series of the continuous function $f$ diverges at $0$. In 1966 Carleson proved that this can happen only on a set of measure zero. Kahane and Katznelson showed that any set of measure zero can occur.
\end{solution}

\begin{solution}{pb:flat}
\idea Parseval says the mean of $|P|^2$ on the circle is the sum of the squares of the coefficients, so the maximum is at least $\sqrt{n+1}$. For an upper bound, choose the signs at random. At each fixed point of the circle, $P$ is a sum of independent bounded terms, so Hoeffding makes a large value exponentially unlikely, and a union bound handles finitely many points. Between the points, $P$ cannot move much, because its derivative is bounded.

\writeup (a) For $z=e^{i\theta}$, one has $|P(z)|^2=\sum_{k,l}\varepsilon_k\varepsilon_le^{i(k-l)\theta}$, and the integral of $e^{im\theta}$ over a period is $0$ for $m\ne0$. So
\[
\frac1{2\pi}\int_0^{2\pi}|P(e^{i\theta})|^2\dd\theta=\sum_k\varepsilon_k^2=n+1,
\]
and $\|P\|^2$, being at least this mean, is at least $n+1$.

(b) Let the $\varepsilon_k$ be independent fair signs. Fix $\theta$. The variables $\varepsilon_k\cos k\theta$ are independent and centred, with values in $[-|\cos k\theta|,|\cos k\theta|]$. Hoeffding's inequality (Volume~4, Chapter~4) gives
\[
P\bigl(|\operatorname{Re}P(e^{i\theta})|\ge t\bigr)\le2\exp\Bigl(-\frac{2t^2}{\sum_k4\cos^2k\theta}\Bigr)\le2e^{-t^2/(2(n+1))},
\]
and the same bound holds for the imaginary part. Let $M=n(n+1)$ and $\theta_j=\frac{2\pi j}M$, for $0\le j<M$. If $|\operatorname{Re}P|<t$ and $|\operatorname{Im}P|<t$, then $|P|<\sqrt2\,t$, so
\[
P\bigl(\exists j:|P(e^{i\theta_j})|\ge\sqrt2\,t\bigr)\le4Me^{-t^2/(2(n+1))}.
\]
With $t^2=2(n+1)\log(5M)$ the right-hand side is $\frac45<1$, so some choice of signs has $|P(e^{i\theta_j})|<\sqrt2\,t=2\sqrt{(n+1)\log(5M)}$ for every $j$. For this choice and any $\theta$, let $\theta_j$ be the nearest grid point, so $|\theta-\theta_j|\le\frac\pi M$. Since $|P'(z)|\le\sum_kk=\frac{n(n+1)}2$ on the circle,
\[
|P(e^{i\theta})-P(e^{i\theta_j})|\le|\theta-\theta_j|\cdot\frac{n(n+1)}2\le\frac\pi M\cdot\frac M2=\frac\pi2.
\]
So $\|P\|\le2\sqrt{(n+1)\log\bigl(5n(n+1)\bigr)}+\frac\pi2$.
\tools Volume~4, Chapter~4 (Hoeffding's inequality, union bound); this chapter (orthogonality of the $e_n$); Volume~1, Chapter~5 (complex numbers).
\leads Salem and Zygmund (1954) showed that for random signs $\|P\|$ is of order $\sqrt{n\log n}$, so (b) is sharp for random
choices; Problem~\ref{pb:rudin} removes the logarithm with a clever choice. Littlewood asked in 1966 for signs with
$\delta\sqrt n\le|P(z)|\le\Delta\sqrt n$ on the whole circle; Balister, Bollob\'as, Morris, Sahasrabudhe and Tiba proved in 2020
that such flat Littlewood polynomials exist.
\end{solution}

\begin{solution}{pb:rudin}
\idea The recursion is built so that $(P_{k+1},Q_{k+1})$ is $(P_k+w,P_k-w)$ with $|w|=|Q_k|$. The parallelogram law then doubles $|P|^2+|Q|^2$ at every point of the circle. The coefficients concatenate. The first cases are $P_2=1+z+z^2-z^3$ and $P_3=1+z+z^2-z^3+z^4+z^5-z^6+z^7$.

\writeup (a) By induction, $P_k$ and $Q_k$ have degree less than $2^k$, and $z^{2^k}Q_k$ has only terms of degree between $2^k$ and $2^{k+1}-1$. The coefficient list of $P_{k+1}$ is that of $P_k$ followed by that of $Q_k$, and the list of $Q_{k+1}$ is that of $P_k$ followed by that of $-Q_k$. Every coefficient is $\pm1$, and the top coefficient is non-zero.

(b) For $|z|=1$, let $u=P_k(z)$ and $v=z^{2^k}Q_k(z)$, so $|v|=|Q_k(z)|$. Then
\[
|P_{k+1}|^2+|Q_{k+1}|^2=|u+v|^2+|u-v|^2=2|u|^2+2|v|^2=2\bigl(|P_k|^2+|Q_k|^2\bigr),
\]
and $|P_0|^2+|Q_0|^2=2$. The sum is $2^{k+1}$, and $|P_k(z)|\le\sqrt{2^{k+1}}$. Problem~\ref{pb:flat}(a) with $n+1=2^k$ gives $\|P_k\|\ge\sqrt{2^k}$.

(c) For $m<2^k$ let $t_k(m)$ be the binary digit of $m$ in position $k-1$, the top position. We prove by induction on $k$ that for $m<2^k$,
\[
[z^m]P_k=(-1)^{e(m)},\qquad[z^m]Q_k=(-1)^{e(m)+t_k(m)}.
\]
For $k=0$ one has $m=0$, and both coefficients are $1$. Assume the formulas for $k$, and let $m<2^{k+1}$. If $m<2^k$, then $[z^m]P_{k+1}=[z^m]Q_{k+1}=[z^m]P_k=(-1)^{e(m)}$, and $t_{k+1}(m)=0$. If $m=2^k+m'$ with $m'<2^k$, then $e(m)=e(m')+t_k(m')$, because the new top digit $1$ forms a new pair exactly when the digit below it is $1$, and $t_{k+1}(m)=1$. So
\[
[z^m]P_{k+1}=[z^{m'}]Q_k=(-1)^{e(m')+t_k(m')}=(-1)^{e(m)}
\]
and $[z^m]Q_{k+1}=-[z^{m'}]Q_k=(-1)^{e(m)+1}$.
\tools Volume~1, Chapter~5 (modulus of complex numbers); this chapter (through Problem~\ref{pb:flat}).
\leads Golay invented such complementary pairs in 1951 for infrared spectrometry. They are used in radar and in Wi-Fi signals. The Rudin--Shapiro sequence $(-1)^{e(m)}$ is $2$-automatic, and Mauduit and Rivat (2015) proved that it takes the values $\pm1$ equally often on the primes.
\end{solution}

\begin{solution}{pb:siegel}
\idea If $K$ avoids the lattice, the half-body $\frac12K$ packs. The periodic function $\varphi$ is then an indicator. Its mean square equals its mean, $2^{-n}\operatorname{vol}K$, while Bessel says the mean square is at least the square of the mean. This forces $2^{-n}\operatorname{vol}K\le1$. The other Fourier coefficients measure how far the packing is from filling space.

\writeup (a) If $x+m=y+m'$ with $x,y\in\frac12K$ and $m\ne m'$, then $m-m'=y-x=\frac12(2y)+\frac12(-2x)$, with $2y\in K$ and $-2x\in K$ by symmetry. Convexity puts $m-m'$ in $K$. That point is a non-zero integer point of $K$, which was excluded.

(b) The set $K$ is bounded, so on $[0,1[^n$ only finitely many terms of $\varphi$ are not identically $0$. By (a), at most one of them is non-zero at each point, so $\varphi\in\{0,1\}$ and $\varphi^2=\varphi$. Unfolding, with the admitted additivity and translation invariance,
\[
\hat\varphi(k)=\sum_m\int_{[0,1[^n}1_{\frac12K}(x+m)e(-k\cdot x)\dd x=\sum_m\int_{[0,1[^n+m}1_{\frac12K}(y)e(-k\cdot y)\dd y=\int_{\frac12K}e(-k\cdot y)\dd y,
\]
since $e(k\cdot m)=1$ and the cubes $[0,1[^n+m$ tile $\R^n$. For $k=0$ this is $\operatorname{vol}\frac12K=2^{-n}\operatorname{vol}K$.

(c) The functions $e(k\cdot x)$ are orthonormal on the torus. The integrals factor into one-dimensional integrals. For a finite set $F\subset\Z^n$,
\[
0\le\int\Bigl|\varphi-\sum_{k\in F}\hat\varphi(k)e(k\cdot x)\Bigr|^2=\int\varphi^2-\sum_{k\in F}|\hat\varphi(k)|^2.
\]
Since $\int\varphi^2=\int\varphi=v$ with $v=2^{-n}\operatorname{vol}K$, one gets $v-v^2\ge\sum_{k\in F\setminus\{0\}}|\hat\varphi(k)|^2$ for every finite $F$. In particular $v(1-v)\ge0$, so $\operatorname{vol}K\le2^n$. A symmetric convex body of volume greater than $2^n$ therefore contains a non-zero integer point.

(d) The open cube contains no non-zero integer point, and $v=1$. For $k\ne0$ some coordinate $k_j$ is non-zero, and
\[
\hat\varphi(k)=\prod_j\int_{-1/2}^{1/2}e(-k_jx_j)\dd x_j=0,
\]
because the factor $j$ vanishes. Both sides are $0$. The translates of $\frac12K$ tile space.
\tools{This chapter (orthogonality, the proof of Bessel's inequality, the section on several variables); the admitted facts on integrals over convex sets.}
\leads With Parseval in place of Bessel, the inequality becomes Siegel's identity $v(1-v)=\sum_{k\ne0}|\hat\varphi(k)|^2$, which measures the defect $2^n-\operatorname{vol}K$ of a lattice-free body. Siegel published it in 1935. Ten years later his mean-value theorem over the space of all lattices opened a chapter of the geometry of numbers, the subject of Volume~13, Chapter~5. The case of a rectangle in the plane was Exercise~1.15 of Volume~2.
\end{solution}

\begin{solution}{ex:rw06a}
\idea At $x=\frac\pi2$ the sines of the odd multiples are $\pm1$, so each partial sum is $\frac4\pi$ times a piece of Leibniz's series.
\writeup One has $\sin\bigl((2k-1)\frac\pi2\bigr)=(-1)^{k-1}$. Therefore
\[
S_1f\Bigl(\frac\pi2\Bigr)=\frac4\pi,\qquad
S_3f\Bigl(\frac\pi2\Bigr)=\frac4\pi\Bigl(1-\frac13\Bigr)=\frac8{3\pi},\qquad
S_5f\Bigl(\frac\pi2\Bigr)=\frac4\pi\Bigl(1-\frac13+\frac15\Bigr)=\frac{52}{15\pi}.
\]
Numerically, $\frac4\pi\approx1.27324$, $\frac8{3\pi}\approx0.84883$, and $\frac{52}{15\pi}\approx1.10347$. The values $S_1f\bigl(\frac\pi2\bigr)$ and $S_5f\bigl(\frac\pi2\bigr)$ lie above $1$. The value $S_3f\bigl(\frac\pi2\bigr)$ lies below $1$. The target of the series at this point is $1$.
\end{solution}

\begin{solution}{ex:rw06b}
\idea The energy of the sawtooth is an elementary integral. Parseval turns it into $\zeta(2)$, and the first five terms are a fraction you can add by hand.
\writeup (a) Substitute $u=\pi-x$. As $x$ runs from $0$ to $2\pi$, the variable $u$ runs from $\pi$ down to $-\pi$, and
\[
\frac1{2\pi}\int_0^{2\pi}\Bigl(\frac{\pi-x}2\Bigr)^2\dd x=\frac1{2\pi}\int_{-\pi}^\pi\frac{u^2}4\dd u=\frac1{2\pi}\cdot\frac12\cdot\frac{\pi^3}3=\frac{\pi^2}{12}.
\]
(b) One has $|c_n|^2=\frac1{4n^2}$ for $n\ne0$, so Parseval reads $\frac{\pi^2}{12}=\sum_{n\ne0}\frac1{4n^2}=\sum_{n\ge1}\frac1{2n^2}$. Since $\zeta(2)=\frac{\pi^2}6$, the same sum is $\frac12\cdot\frac{\pi^2}6=\frac{\pi^2}{12}$.
(c) The sum of the first five reciprocals of squares is
\[
1+\frac14+\frac19+\frac1{16}+\frac1{25}=\frac{3600+900+400+225+144}{3600}=\frac{5269}{3600},
\]
so $\sum_{n=1}^5\frac1{2n^2}=\frac{5269}{7200}\approx0.73181$. And $\frac{\pi^2}{12}\approx0.82247$. The gap is about $0.09066$. The first five harmonics carry most of the energy, and the tail is still visible.
\end{solution}

\begin{solution}{ex:rw06c}
\idea Rationality of $\alpha$ would make $\sqrt5$ rational. The count for $n\alpha$ is a sample of Corollary~\ref{cor:nalpha}. The count for $n^2/5$ is a count of quadratic residues.
\writeup (a) The number $\alpha$ satisfies $\alpha^2+\alpha-1=0$, so $2\alpha+1=\sqrt5$. If $\alpha=\frac pq$, then $\sqrt5$ would be rational. The usual argument for $\sqrt2$ applies. If $\sqrt5=\frac ab$ in lowest terms, then $a^2=5b^2$, so $5$ divides $a$. Writing $a=5c$ gives $5c^2=b^2$, so $5$ divides $b$, a contradiction.
(b) The twenty fractional parts, to three decimals, are
\[
\begin{aligned}
&0.618,\ 0.236,\ 0.854,\ 0.472,\ 0.090,\ 0.708,\ 0.326,\ 0.944,\\
&0.562,\ 0.180,\ 0.798,\ 0.416,\ 0.034,\ 0.652,\ 0.271,\ 0.889,\\
&0.507,\ 0.125,\ 0.743,\ 0.361.
\end{aligned}
\]
Exactly ten of them lie in $\bigl[0,\frac12\bigr)$, namely those of $n=2,4,5,7,10,12,13,15,18,20$. Ten out of twenty is the share the interval should get.
(c) Modulo $5$, the square $n^2$ is $0$, $1$, or $4$, according as $n\equiv0$, $n\equiv\pm1$, or $n\equiv\pm2$. So $\{n^2/5\}$ is $0$, $\frac15$, or $\frac45$. The first two lie in $\bigl[0,\frac12\bigr)$, and $\frac45$ does not. Among each complete set of residues there are three favourable classes out of five, and $n=1,\ldots,20$ contains four of each residue. The count is $12$. The proportion of $n\le N$ with $\{n^2/5\}\in\bigl[0,\frac12\bigr)$ tends to $\frac35$, because the residues of $n$ modulo $5$ are equally frequent and three of the five squares land in the interval.
\end{solution}

\begin{solution}{ex:rw06d}
\idea An irrational step is equidistributed by the corollary. A rational step is a finite set of points, and one interval can be chosen to hold the wrong proportion of them. The geometric sum in the proof of the corollary is a number you can compute for $N=10$.
\writeup (a) If $\sqrt2=\frac ab$ in lowest terms, then $a^2=2b^2$, so $a$ is even. Writing $a=2c$ gives $2c^2=b^2$, so $b$ is even, a contradiction. Thus $\sqrt2$ is irrational, and Corollary~\ref{cor:nalpha} says that $(n\sqrt2)$ is equidistributed modulo $1$.
(b) Let $q=e^{2\pi i\sqrt2}$. Then $q\ne1$, and
\[
\sum_{n=1}^{10}q^n=\frac{q(q^{10}-1)}{q-1}.
\]
The modulus of the left side is $0.448$. The bound in the proof of the corollary is $\frac2{|q-1|}=1.037$. The inequality holds, and the average of the ten terms is $0.0448$, already small.
(c) The fractional parts $\{n/3\}$ cycle through $\frac13$, $\frac23$, $0$. The proportion in $\bigl[0,\frac12\bigr)$ tends to $\frac23$, not to $\frac12$, so $(n/3)$ is not equidistributed. Among $n=1,\ldots,30$, the inequality $\{n\sqrt2\}<0.1$ holds for $n=5$, $n=17$ and $n=29$, three times, and $30\times0.1=3$. In the same range, $\{n/3\}$ lies in $\bigl[0,\frac12\bigr)$ for $20$ indices, which is exactly $\frac23$ of $30$.
\end{solution}

\section*{Sources and further reading}

\begin{itemize}
\item T.~W.~K\"orner, \emph{Fourier Analysis}, Cambridge University Press, 1988. Short chapters, full of history and
applications; the Gibbs phenomenon, Weyl and the heat equation are all there.
\item E.~Stein and R.~Shakarchi, \emph{Fourier Analysis: An Introduction}, Princeton University Press, 2003, Chapters~2--4
(including Hurwitz's isoperimetric proof and Weyl's criterion).
\item Y.~Katznelson, \emph{An Introduction to Harmonic Analysis}, 3rd ed., Cambridge University Press, 2004.
\item A.~Hurwitz, \emph{Sur quelques applications g\'eom\'etriques des s\'eries de Fourier}, Annales scientifiques de l'\'E.N.S.
19 (1902), 357--408.
\item H.~Weyl, \emph{\"Uber die Gleichverteilung von Zahlen mod.~Eins}, Mathematische Annalen 77 (1916), 313--352.
\item J.~Fourier, \emph{Th\'eorie analytique de la chaleur}, Firmin Didot, Paris, 1822. The sentence in the opening historical paragraph is from the
\emph{Discours pr\'eliminaire}. For the memoir of 1807 and Lagrange's objection: I.~Grattan-Guinness, \emph{Joseph Fourier
1768--1830}, MIT Press, 1972.
\item L.~Fej\'er, \emph{Sur les fonctions born\'ees et int\'egrables}, Comptes Rendus de l'Acad\'emie des Sciences 131 (1900),
984--987, for Theorem~\ref{thm:fejer}; the construction of Exercise~\ref{ex:divergent} follows his later papers.
\item L.~Carleson, \emph{On convergence and growth of partial sums of Fourier series}, Acta Mathematica 116 (1966), 135--157;
J.-P.~Kahane and Y.~Katznelson, \emph{Sur les ensembles de divergence des s\'eries trigonom\'etriques}, Studia Mathematica 26
(1966), 305--306.
\item W.~Rudin, \emph{Some theorems on Fourier coefficients}, Proceedings of the American Mathematical Society 10 (1959),
855--859; P.~Balister, B.~Bollob\'as, R.~Morris, J.~Sahasrabudhe and M.~Tiba, \emph{Flat Littlewood polynomials exist},
Annals of Mathematics 192 (2020), 977--1004; C.~Mauduit and J.~Rivat, \emph{Prime numbers along Rudin--Shapiro sequences},
Journal of the European Mathematical Society 17 (2015), 2595--2642.
\item C.~L.~Siegel, \emph{\"Uber Gitterpunkte in convexen K\"orpern und ein damit zusammenh\"angendes Extremalproblem}, Acta
Mathematica 65 (1935), 307--323, for Exercise~\ref{pb:siegel}.
\item E.~Hewitt and R.~Hewitt, \emph{The Gibbs--Wilbraham phenomenon: an episode in Fourier analysis}, Archive for History of
Exact Sciences 21 (1979), 129--160.
\end{itemize}
